%%%%%%%%%%%%%%%%%%%%%%%%%%%%%%%%%%%%%%%%%%%%%%%%%%%%%%
\chapter{General conclusions}
\label{general_conclusions}
\graphicspath{{chapter_09/figures}{chapter_09/tables}}
%%%%%%%%%%%%%%%%%%%%%%%%%%%%%%%%%%%%%%%%%%%%%%%%%%%%%%

Flash floods claim over 5,000 lives annually and represent the most devastating form of flood hazard worldwide, with their impacts transcending traditional divides between Global North and South. Recent catastrophic events in the USA, Spain, Libya, Germany/Belgium, and China have underscored the urgent need for enhanced early warning capabilities, particularly as climate change intensifies extreme rainfall events, even in historically low-risk regions. This thesis has addressed a critical gap in global disaster preparedness by developing and demonstrating a proof-of-concept for medium-range predictions of areas at risk of flash floods over a continuous global domain. 

The research presented in this thesis makes three fundamental contributions to flash flood science and operational forecasting. First, through the development of a flash-flood-focused verification framework (Chapter 5), this work establishes that global numerical weather prediction models can identify areas at risk of flash floods with meaningful skill up to medium-range timescales (5 days ahead), challenging the prevailing assumption that flash flood prediction must use high-resolution forecasts to generate useful predictions. Second, the successful implementation of data-driven hydro-meteorological models (Chapter 6) demonstrates that machine learning approaches can extract predictive signals from severely imbalanced datasets, in which flash flood events are reported less than 0.3\% of the time. Third, and most significantly for global applications, this research proves that models trained exclusively on high-quality observations from data-rich regions can successfully predict flash flood risk in entirely different geographical contexts (Chapter 7).

The methodology developed here offers a pragmatic pathway toward this ambitious goal, which has particular relevance for the billions of people residing in data-scarce regions where traditional forecasting approaches remain economically or technically unfeasible. By leveraging globally available global numerical weather prediction outputs from systems like ERA5 and employing transfer learning principles, this approach circumvents the fundamental obstacle that has historically prevented global flash flood warning coverage: the absence of comprehensive observational networks in most parts of the world.

The operational implications of this research extend beyond technical achievements. Emergency managers could henceforth access flash flood risk assessments with sufficient lead time to implement graduated response protocols. Day-one forecasts, maintaining performance metrics comparable to reanalysis-based predictions, support high-confidence decisions including evacuations and resource mobilisation. Medium-range forecasts at days three to five, whilst experiencing some skill degradation, retain sufficient accuracy for preparedness activities such as pre-positioning emergency supplies, issuing public advisories, and activating monitoring protocols. This temporal stratification of forecast utility enables cost-effective risk management strategies that balance the imperative to protect lives against the economic and social costs of false alarms.

Several limitations constrain the immediate operational deployment of this system and merit acknowledgement. The 31-kilometre spatial resolution inherited from ERA5 cannot fully capture the localised processes generating flash floods in small catchments, and so cannot pinpoint elevated risks that might be intrinsic to narrow valleys or urban areas. The reliance on impact databases for model training introduces systematic biases related to population density and socio-economic factors that may not transfer uniformly to all global contexts. The absence of comprehensive global verification data prevents rigorous statistical validation of model performance in most regions, necessitating continued reliance on case-study validation. These constraints, whilst significant, do not negate the fundamental advance this research represents in extending flash flood warnings to previously unprotected populations.

Future research should pursue four strategic directions to enhance global flash flood prediction capabilities. First, integration with higher-resolution atmospheric models, including convection-permitting implementations, could substantially improve the representation of localised extreme rainfall events, most notably for shorter lead times. Second, ensemble approaches incorporating multiple data-driven architectures has potential to provide better uncertainty quantification for risk-based decision-making. Third, systematic aggregation of high-density, regional flash flood databases could create more globally representative training datasets. Fourth, the development of causal inference techniques could address spurious correlations identified through machine learning, ensuring predictions align with physical understanding whilst maintaining statistical accuracy.

The vision of universal flash flood protection, advocated by initiatives such as the UN's "Early Warnings for All" initiative, appeared technically unattainable when announced in 2022. This thesis demonstrates that the integration of global numerical weather prediction models, data-driven methodologies, and transfer learning can transform this vision into operational reality. The successful prediction of devastating flash floods in Valencia and Zhengzhou using models trained exclusively on North American data supports the fundamental hypothesis that hydro-meteorological relationships governing flash flood generation exhibit sufficient commonality to enable knowledge transfer across continents. As climate change drives communities into unfamiliar hazard regimes, this capability to extend protective warnings beyond traditional observational boundaries becomes not merely advantageous but essential for global resilience. The proof-of-concept established through this research provides the scientific foundation for a new generation of early warning systems that can help protect vulnerable communities worldwide, ultimately contributing to the reduction of flash flood mortality and the enhancement of global disaster preparedness.